% Source: Wald GR, physical PDF pp. 42--43
%         (printed pp. 29--30).
% Coverage: Chapter 3 opening material before Section 3.1.
% Figures/tables/footnotes: Figures 3.1--3.2; no footnotes.
% Status: source-reviewed and compile-clean.
% Uncertainties: none.

Our intuitive notion of curvature arises mainly from two-dimensional surfaces which
are embedded in ordinary three-dimensional Euclidean space. We normally think of
a surface as curved because of the way it bends in $\mathbb R^3$. In chapter 9, we will capture
this notion by defining the \emph{extrinsic} curvature of a surface embedded in a higher
dimensional space. However, our interest here is to investigate the curvature of
spacetime. Our spacetime manifold $M$ with spacetime metric $g_{ab}$ is not naturally
embedded (at least so far as we know) in a higher dimensional space. Thus we wish
to develop an \emph{intrinsic} notion of curvature that can be applied to any manifold
without reference to a higher dimensional space in which it might be embedded.

Such a notion of curvature can be defined in terms of parallel transport. On a
surface such as a plane (Fig.~\ref{fig:3.1}) or sphere (Fig.~\ref{fig:3.2}), we have an intuitive notion
(which will be made mathematically precise below) of what it means to keep a vector
\lq\lq pointing in the same direction\rq\rq\ (but always in the tangent space of the manifold) as
one moves it along a path. On the plane, if one parallel-transports a vector around
any closed path, the final vector always coincides with its initial value. However, this
is not so on the sphere. The vector shown in Fig.~\ref{fig:3.2} comes back rotated with
respect to its initial value when carried along the path shown. This basic idea allows
us to characterize the plane as flat, the sphere as curved, and more generally allows
us to characterize the curvature of any manifold intrinsically once we are told how
to \lq\lq parallel transport\rq\rq\ vectors along curves.

An alternative characterization of curvature also can be given as follows. A
geodesic is a curve whose tangent is parallel-transported along itself, that is, it is a
\lq\lq straightest possible\rq\rq\ curve. A space will be curved if and only if some initially
parallel geodesics fail to remain parallel, i.e., Euclid's fifth postulate fails.

Given only the manifold structure of space, we do not have a natural notion of
parallel transport. The reason is that the tangent spaces $T_pM$ and $T_qM$ of two distinct points
$p$ and $q$ are \emph{different} vector spaces and hence there is no way of saying that a vector

% Source: Wald GR, physical PDF p. 42 (printed p. 29).
% Coverage: Figure 3.1.
% Figures/tables/footnotes: Figure only; no tables or footnotes.
% Status: source-reviewed and compile-clean.
% Uncertainties: none.
\begingroup
\renewcommand{\thefigure}{3.1}
\begin{figure}[H]
  \centering
  \begin{tikzpicture}[scale=0.92,>=Stealth]
    \draw[thick] (-2.7,-0.8) -- (-1.0,0.8) -- (2.7,0.8) -- (1.2,-0.8) -- cycle;
    \draw[thick] plot[smooth cycle,tension=.75] coordinates
      {(-1.25,-0.1) (-.55,.40) (.82,.37) (1.25,-.08) (.56,-.37) (-.52,-.33)};
    \fill (-1.25,-0.1) circle (1.6pt) node[below left=1pt] {$p$};
    % The transported vector stays parallel to itself: every copy points
    % in the same (horizontal) direction, which is the point of the figure.
    \foreach \x/\y in {-.60/.42,.55/.40,1.28/-.06,.30/-.42,-.70/-.36} {
      \draw[->,thick] (\x,\y) -- ++(.42,0);
    }
    \draw[->,thick] (-1.20,-.10) -- ++(.42,0);
    \node[below=1pt] at (-.99,-.14) {$v^a$};
  \end{tikzpicture}
  \caption{The parallel transport of a vector, $v^a$, around a closed curve in the plane.
  The vector $v^a$ always ``comes back'' pointing in the same direction as it did initially.}
  \label{fig:3.1}
\end{figure}
\endgroup

% Source: Wald GR, physical PDF p. 43 (printed p. 30).
% Coverage: Figure 3.2.
% Figures/tables/footnotes: Figure only; no tables or footnotes.
% Status: source-reviewed and compile-clean.
% Uncertainties: none.
\begingroup
\renewcommand{\thefigure}{3.2}
\begin{figure}[htbp]
  \centering
  \begin{tikzpicture}[scale=.82,>=Stealth]
    \draw[thick] (0,0) circle (1.8);
    % Three mutually orthogonal great-circle arcs: p -> R along the "equator",
    % R -> A up the right arc, A -> p back down the left arc.
    \path[draw,thick] (-1.05,-.50) .. controls (-.30,-.95) and (.75,-.70) .. (1.40,-.12);
    \path[draw,thick] (1.40,-.12) .. controls (1.35,.60) and (.75,1.15) .. (.05,1.55);
    \path[draw,thick] (.05,1.55) .. controls (-.55,1.25) and (-.95,.30) .. (-1.05,-.50);
    % Copies of the vector during transport.  At p show both the initial
    % vector (along the lower arc) and the returned vector (up the left arc):
    % these are separated by 90 degrees in the tangent plane.
    \draw[->,thick] (-.38,-.70) -- ++(.43,.02);
    \draw[->,thick] (.62,-.50) -- ++(.42,.02);
    \draw[->,thick] (.70,1.05) -- ++(.29,.30);
    \draw[->,thick] (1.03,.64) -- ++(.34,.25);
    \draw[->,thick] (.05,1.55) -- ++(.24,.35);
    \draw[->,thick] (-.47,1.08) -- ++(.24,.35);
    \draw[->,thick] (-.85,.30) -- ++(.14,.40);
    \fill (-1.05,-.50) circle (1.6pt) node[below left=1pt] {$p$};
    \draw[->,thick] (-1.02,-.50) -- ++(.44,0);
    \draw[->,thick] (-1.02,-.47) -- ++(.15,.42);
    \node[below=1pt] at (-.80,-.54) {$v^a$};
    % Reserve trim clearance above the artwork when this float lands at page top.
    \path[draw=none] (0,3.00) circle[radius=.1pt];
  \end{tikzpicture}
  \caption{The parallel transport of a vector, $v^a$, around a closed curve on the
  sphere. In the case shown here of a closed curve composed of three mutually
  orthogonal segments of great circles, the vector $v^a$ comes back rotated by $90^{\circ}$.}
  \label{fig:3.2}
\end{figure}
\endgroup


at $p$ is the same as a vector at $q$. Thus, the definition of parallel transport requires
more than just the manifold structure. It is not difficult to convince oneself that a
notion of how to parallel-transport vectors should be equivalent to the knowledge of
how to take derivatives of vector fields. If we know how to parallel-transport vectors
along a curve, we can define the derivative of a vector field in the direction of the
curve; similarly, given a notion of derivative, we can define a vector to be parallel
transported if its derivative along the given curve is zero. It turns out to be most
convenient to work directly with the notion of a derivative operator, and we shall do
so in this chapter. The failure of a vector to return to its original value when parallel
transported around an infinitesimal closed curve translates into the lack of commutativity of derivatives. Thus, the notion of curvature can be defined in terms of the
failure of successive differentiations on tensor fields to commute. This is the route
we shall follow in section 3.2.

From where does this extra structure needed to define parallel transport or a
derivative operator arise? We will show in section 3.1 that given a metric (of any
signature), there is a unique definition of parallel transport which preserves inner
products of all pairs of vectors. Thus, the existence of a metric gives rise to a unique
notion of parallel transport and, thus, to an intrinsic notion of the curvature of the
manifold. This is the notion of the curvature of a spacetime $(M,g_{ab})$ in which we are
interested. However, it is more convenient to proceed by giving a general definition
of the notions of derivative operator, parallel transport, and curvature before specializing to the case where they arise from a metric, and we shall proceed in this manner.

Derivative operators and parallel transport are defined in section 3.1, and curvature is defined in section 3.2. In much of our discussion in these sections, we shall
follow closely the treatment given in unpublished notes of Geroch. Geodesics are
introduced in section 3.3, and the geodesic deviation equation---which characterizes
curvature in terms of the failure of initially parallel geodesics to remain parallel---is
derived. Finally, section 3.4 discusses methods for computing curvature.

